\subsection{Conditional Independence}
\totalscore{9}

Let $\Xcal,\Ycal,\Zcal,\Tcal$ be standard measurable spaces.
Consider three conditional random variables $X,Y,Z$ with codomains $\Xcal$, $\Ycal$, $\Zcal$, respectively.
Their joint Markov kernel is denoted by $K(X,Y,Z|T)$.
%If you like, you are allowed to assume (without further justification) that \emph{all occurring Markov kernels are discrete and have mass functions}.
%(without loss of points). 
%and the corresponding mass function by $k$, in arguments written as $k(x,y,z|t)$.

Recall that $X$ is independent from $Y$ conditioned on $Z$ w.r.t.\ $K(X,Y,Z|T)$, in symbols:
\begin{align}
    X \Indep_{K(X,Y,Z|T)} Y \given Z,
\end{align}
iff there exists a Markov kernel $Q(X|Z) :\, \Zcal \dshto \Xcal$ such that:
\begin{align}
    K(X,Y,Z|T) &= Q(X|Z) \otimes K(Y,Z|T).
\end{align}

Now let $P(T)$ be any fixed probability distribution on $\Tcal$.
We then abbreviate the joint probability distribution of $X,Y,Z$ and $T$ as follows:
\begin{align}
    P(X,Y,Z,T) &:= K(X,Y,Z|T) \otimes P(T).
\end{align}


\begin{enumerate}[label=\alph*)]
    \item Show that we always have the implication: \score{3}
      \begin{align}
            X \Indep_{K(X,Y,Z|T)} Y \given Z \qquad \implies \qquad X \Indep_{P(X,Y,Z,T)} Y,T \given Z.
      \end{align}
  \item  Now assume that $\Tcal$ is a countable discrete space and that $P(T)$ has a strictly positive mass function $p$, i.e.\ for all $t \in \Tcal$ we have: \score{6}
\begin{align}
    p(t) & >0.
\end{align}
        Show that in this case we also have the reverse implication:
      \begin{align}
            X \Indep_{K(X,Y,Z|T)} Y \given Z \qquad \impliedby \qquad X \Indep_{P(X,Y,Z,T)} Y,T \given Z.
      \end{align}
\end{enumerate}   
Clearly exhibit all arguments involved in your reasoning. 


\answer{
``$\implies$'': By definition of $X \Indep_{K(X,Y,Z|T)} Y \given Z$ we have a factorization:
\begin{align}
    K(X,Y,Z|T) &= Q(X|Z) \otimes K(Y,Z|T),
\end{align}
where $K(Y,Z|T)$ is the marginal of $K(X,Y,Z|T)$. Multiplying each side with $P(T)$ gives us:
\begin{align}
    P(X,Y,Z,T) & = K(X,Y,Z|T) \otimes P(T) \\
               &= Q(X|Z) \otimes K(Y,Z|T) \otimes P(T) \\
               &= Q(X|Z) \otimes P(Y,Z,T).
\end{align}
This shows the defining property of $X \Indep_{P(X,Y,Z,T)} Y,T \given Z$, where $Q(X|Z)$ is the wanted Markov kernel.

``$\impliedby$'': Assume $X \Indep_{P(X,Y,Z,T)} Y,T \given Z$. Then there exists a Markov kernel $Q(X|Z)$ such that:
\begin{align}
    P(X,Y,Z,T) &= Q(X|Z) \otimes P(Y,Z,T),
\end{align}
where $P(Y,Z,T)$ is the marginal of $P(X,Y,Z,T)$. 
This shows:
\begin{align}
    K(X,Y,Z|T) \otimes P(T) 
               & = P(X,Y,Z,T) \\
               &= Q(X|Z) \otimes P(Y,Z,T) \\
               &= Q(X|Z) \otimes K(Y,Z|T) \otimes P(T).
\end{align}
Again note that for the marginals we also have:
\begin{align} 
    P(Y,Z,T) = K(Y,Z|T) \otimes P(T).
\end{align}
By the essential uniqueness of conditional Markov kernels we get:
\begin{align}
    K(X,Y,Z|T) &= Q(X|Z) \otimes K(Y,Z|T)  \qquad P(T)\text{-almost-surely.}
\end{align}
More precisely, this means that the following measurable set:
\begin{align}
    N &:= \lC t \in \Tcal \st K(X,Y,Z|T=t) \neq Q(X|Z) \otimes K(Y,Z|T=t)   \rC
\end{align}
is a $P(T)$-null set. Since $P(T)$ is strictly positive the set $N$ must be empty.
So for all $t \in \Tcal$ we have:
\begin{align}
    K(X,Y,Z|T=t) &= Q(X|Z) \otimes K(Y,Z|T=t),
\end{align}
which implies:
\begin{align}
    K(X,Y,Z|T) &= Q(X|Z) \otimes K(Y,Z|T),
\end{align}
and thus: $X \Indep_{K(X,Y,Z|T)} Y \given Z$.
}
